\documentclass[
    12pt,
    a4paper,
    sumario=tradicional,
    oneside,
    article
]{abntex2}

% ---
% Pacotes básicos
% ---
\usepackage[utf8]{inputenc}
\usepackage[T1]{fontenc}
\usepackage{lmodern}
\usepackage[brazil]{babel}
\usepackage{graphicx}
\usepackage{indentfirst}
\usepackage[alf]{abntex2cite}

% ---
% Dados do trabalho
% ---
\titulo{Uma Análise Empírica entre Categorias ODC e Eficácia de Testes via Mutation Testing}
\autor{Guilherme de Almeida Rocha Vieira \and Felipe Augusto Pereira de Sousa}
\local{Belo Horizonte}
\data{2026}
\instituicao{Pontifícia Universidade Católica de Minas Gerais}

\begin{document}

% ---
% Elementos pré-textuais
% ---
\pretextual

\imprimircapa

% Resumo em português
\begin{resumo}
A qualidade dos testes de software é fundamental para garantir a
confiabilidade de sistemas, mas métricas tradicionais, como a
cobertura de código, são insuficientes para capturar se os testes
de fato detectam defeitos. Além disso, a ausência de uma
caracterização do tipo de defeito que escapa aos testes dificulta
o direcionamento de esforços de melhoria da estratégia de teste de
forma orientada a dados. Este trabalho investiga a relação entre
categorias da Classificação Ortogonal de Defeitos (ODC) e a
eficácia de testes, medida por meio de mutation testing, buscando
verificar se diferentes tipos de defeito apresentam níveis
distintos de detecção. Para isso, pull requests de correção de
bugs foram minerados de repositórios open-source no GitHub,
classificados segundo os atributos do ODC e avaliados por meio de
mutation testing (via Stryker), comparando a capacidade de
detecção da suíte de testes antes e depois da correção do defeito.
[RESULTADO A PREENCHER]. Este estudo contribui para a literatura
de engenharia de software ao integrar uma taxonomia consolidada de
causa de defeito com uma métrica de eficácia de teste mais
rigorosa que a cobertura, oferecendo subsídios para estratégias de
teste direcionadas por tipo de defeito.

 \noindent
 \textbf{Palavras-chave}: Teste de software; Mutation testing; ODC; Detecção de defeitos.
\end{resumo}

% Resumo em inglês (abstract)
\begin{resumo}[Abstract]
Software test quality is critical to ensure system reliability, but
traditional metrics, such as code coverage, are insufficient to
capture whether tests actually detect defects. Furthermore, the
lack of characterization of the defect type that escapes testing
makes it difficult to direct test strategy improvement efforts in
a data-driven manner. This work investigates the relationship
between Orthogonal Defect Classification (ODC) categories and test
effectiveness, measured through mutation testing, seeking to
verify whether different defect types exhibit distinct detection
levels. To this end, bug-fixing pull requests were mined from
open-source GitHub repositories, classified according to ODC
attributes, and evaluated through mutation testing (using Stryker),
comparing the test suite's detection capability before and after
the defect fix. [RESULT TO BE FILLED IN]. This study contributes to
the software engineering literature by integrating, for the first
time, a well-established defect-cause taxonomy with a test
effectiveness metric more rigorous than coverage, offering insights
for defect-type-driven, data-oriented test strategies.

 \vspace{1em}
 \noindent
 \textbf{Keywords}: Software testing; Mutation testing; ODC; Defect detection.
\end{resumo}

% ---
% Sumario
% ---
\pdfbookmark[0]{\contentsname}{toc}
\tableofcontents*
\cleardoublepage

% ---
% Elementos textuais
% ---
\textual

\section{Introdução}

A qualidade dos testes de software é uma questão crítica para qualquer
aplicação, porém, métricas tradicionais, tais como a cobertura de testes,
são insuficientes para medir se os testes realmente detectam defeitos
\cite{papadakis2018}.
Um obstáculo comum é a falta de definição do tipo de defeito que escapa
dos testes, fazendo com que equipes não consigam melhorar a estratégia
de testes de forma direcionada e orientada a dados.

Nesse contexto, a taxonomia ODC (Orthogonal Defect Classification)
\cite{chillarege1992} oferece um mecanismo consolidado para classificar
as causas dos defeitos de software, enquanto o mutation testing fornece
uma métrica de eficácia de testes com maior rigor do que a simples
cobertura de linhas de código \cite{papadakis2018}.
A correlação entre essas duas abordagens, no entanto, ainda não foi
investigada empiricamente, configurando uma lacuna relevante na
literatura \cite{butcher2002, chillarege1992}.

Este artigo investiga se categorias ODC distintas apresentam diferentes
níveis de detecção quando avaliadas por meio do mutation score.

% ---
% Elementos pós-textuais
% ---
\postextual

\bibliography{bibliografia/references}

\end{document}